\section{Product Background}

The rapid growth of remote work has turned freelance collaboration into a default mode of software and digital service delivery. On large marketplaces, however, the same problems keep recurring. Clients cannot easily judge whether a candidate is who the profile claims to be, and they must commit money before any deliverable exists. Freelancers, on the other side, cannot easily verify that a client will pay, and they often work without a written scope, without a milestone schedule, and without a recoverable record of what was agreed. The result is a market where most of the effort is spent on risk mitigation rather than on the work itself, and where the tools that would mitigate that risk -- identity verification, escrow, structured milestones, evidence-based dispute handling -- are usually sold as separate products or reserved for enterprise customers.

FREVIA (Freelance Via) is a web-based freelance marketplace that addresses this problem by putting verification, escrow, and project governance into the core transaction rather than into an optional add-on. The platform covers the complete engagement lifecycle: a client publishes a job with a fixed budget, deadline, required skills, and the number of positions to fill; freelancers submit proposals, either as drafts or as submitted applications with attachments and an expiry date; the client compares proposals and hires one candidate; a contract is generated from the accepted proposal, priced at the agreed total amount with a fixed platform fee recorded on the contract, and must be signed by both parties before work starts; the work itself is divided into milestones that are funded, submitted, reviewed, and released one at a time; and the engagement closes with a review that becomes part of both parties' public reputation.

Trust is enforced at several layers. Every freelancer can complete an electronic identity verification in which the front and back of an identity document, a selfie, and a short sequence of liveness frames are submitted to an AI verification service that performs optical character recognition, QR-code reading, image quality assessment, face matching, passive and active liveness detection, and tamper detection before returning a verified or rejected result. Verified documents are additionally reviewed by an administrator, so a freelancer profile can display an identity badge that is backed by evidence rather than by a self-declaration. Profile edits made by freelancers, including expert profiles, are submitted as revisions with a proposed data snapshot and a profile strength score, and only become public after administrative approval. On the financial side, client funds are held in escrow at the milestone level, so a freelancer can see that money exists before starting a milestone, and a client only releases funds after reviewing the submitted deliverable.

Disagreement is treated as a first-class process rather than an exception. When a milestone submission is rejected and the two parties cannot agree, either side can open a dispute on that milestone. The dispute carries an arbitration fee with its own payment deadline and a response deadline, both parties file structured evidence marked as claim or response, and an administrator issues a decision that splits the milestone amount between freelancer and client and records a written reasoning. A party that disagrees with the decision may request a review, after which a final decision is recorded and the dispute is closed. Every resulting money movement is written to a transaction ledger with a distinct type -- escrow deposit, milestone payout, platform fee, dispute fee, arbitrator payout, or refund -- so the financial history of a contract can be reconstructed.

Artificial intelligence is applied in three narrowly scoped places where it removes real manual work, and deliberately not in the place where it would take a decision away from the user. First, a freelancer can upload a PDF curriculum vitae and have an external AI service extract a list of skills together with an estimated proficiency level on a one-to-ten scale; the freelancer reviews the suggestion and decides what to add to the profile. Second, the identity verification described above is an AI pipeline rather than a manual document check. Third, forum content is screened by an AI moderation step that returns a score, matched categories, and a recommendation; every post still enters a pending state and becomes publicly readable only after an administrator approves it. FREVIA does not implement an automated job-to-freelancer recommendation or matching engine, because hiring remains a human judgement that depends on context the platform does not observe, and because a ranking model trained on a young marketplace would encode the biases of a very small sample.

Collaboration during the engagement is supported by a real-time messaging layer built on WebSocket connections authenticated with the same access token as the rest of the API. A conversation is one-to-one between a client and a freelancer, supports text and file attachments, read receipts, typing indicators, per-participant hiding and pinning, and soft deletion of individual messages. Contracts additionally expose shared files, so that attachments do not have to live in a private chat history. A persistent notification record is written for every relevant domain event, and a community forum with categories, comments, likes, and reporting gives users a place for public discussion that is moderated rather than unmoderated.

The platform is bilingual. The interface is available in English and Vietnamese, selected through a locale cookie, and language switching does not change the URL structure, which keeps deep links stable when a user changes language. Access control is dynamic rather than hard-coded: a permission is a record that binds an HTTP method and a path pattern to a module, roles are granted permissions, and a single guard evaluates every request against the permissions of the role currently active in the session. This makes it possible for an administrator to adjust what a role can do without a redeployment, and it allows one account to hold several roles and switch between them, which matters for users who both hire and work.

Financially, FREVIA relies entirely on Stripe. Freelancers and clients connect a Stripe Express account through an onboarding link, and the onboarding status is tracked on the user record so the platform knows whether payouts are possible. Milestone funding can be started either as a PaymentIntent or as a hosted Checkout Session, and a saved card can be reused for later milestones. Released milestones are paid out as Stripe transfers, refunds are issued as Stripe refunds, and all asynchronous state changes arrive through a single webhook endpoint whose signature is verified against the raw request body. No blockchain component, wallet, or proprietary settlement infrastructure is involved; the transaction ledger is an application-level record that mirrors the state reported by Stripe.

FREVIA is therefore best understood as a governed marketplace rather than a listing board. Its value comes from the sequence it enforces -- verified identity, signed contract, funded milestone, reviewed deliverable, recorded decision, rated outcome -- and from the fact that every step of that sequence leaves a durable record that both parties and the platform administrator can inspect later.
